\documentclass[11pt]{article}
\usepackage[T1]{fontenc}
\usepackage{lmodern,amsmath,amssymb,amsthm,microtype}
\usepackage[letterpaper,margin=1in]{geometry}
\usepackage[colorlinks=true,urlcolor=blue,linkcolor=blue]{hyperref}
\hypersetup{pdftitle={Right-half-plane independence roots of random graphs}}
\newtheorem{theorem}{Theorem}
\newtheorem{lemma}[theorem]{Lemma}
\newtheorem{corollary}[theorem]{Corollary}
\newcommand{\E}{\mathbb E}
\newcommand{\Prob}{\mathbb P}
\newcommand{\Unst}{\mathcal U}
\newcommand{\Rea}{\operatorname{Re}}
\newcommand{\pois}{\operatorname{Poisson}}
\title{Right-half-plane independence roots of random graphs}
\author{Research note}
\date{18 November 2026}
\begin{document}
\maketitle
\begin{abstract}
Brown and Cameron asked whether almost all graphs have all their
independence roots in the closed left half-plane. We give a negative
answer for every fixed edge probability $p\in(0,1)$. The limit superior
of the probability of a root in the open right half-plane is at least
$1-\exp(-(1-p)^2)$. In particular it exceeds $22\%$ in the uniform
labeled graph model. This also disproves Cameron's sector conjecture.
The proof tunes the expected number of maximum independent sets to a
constant. One extra maximum set can then violate a necessary Hurwitz
inequality while the next three coefficients concentrate. On suitable
subsequences roots can reach real part of order $n/(\log n)^2$ with
probability bounded away from zero.
\end{abstract}

\section{The question and the obstruction}
Let $G=G(n,p)$ be the labeled random graph with independent edges of
probability $p$. Write $X_s$ for the number of independent $s$-sets,
$\alpha$ for the independence number, and $I_G(z)=\sum_{s=0}^\alpha X_s z^s$.
A graph is stable if all zeros of $I_G$ have nonpositive real part.
Write $\Unst_n$ for the event that $I_G$ has a zero of strictly positive
real part. Brown and Cameron asked whether $\Prob(\Unst_n)\to0$.
Cameron's stronger Conjecture 6.1.1 places all roots of almost every
graph in the sector of angular width $2\pi/3$ about the negative
real axis. That sector lies in the left half-plane.
[\href{https://doi.org/10.37236/7280}{Brown and Cameron, Problem 2};
\href{https://hdl.handle.net/10222/75412}{Cameron, Conjecture 6.1.1 and Problem 6.2.1}]

\begin{theorem}\label{thm:main}
For every fixed $p\in(0,1)$ and $q=1-p$,
\[
 \limsup_{n\to\infty}\Prob(\Unst_n)\ \ge\ 1-e^{-q^2}.
\]
In particular, for $p=1/2$ the bound is $1-e^{-1/4}=0.221199\ldots$.
\end{theorem}

If a polynomial of degree $k$ has positive leading coefficient, denote
its coefficients in descending order by $a_0,a_1,\ldots,a_k$.
A necessary condition for all roots to have nonpositive real part is
\[
 D_2=a_1a_2-a_0a_3\ge0.                                      \tag{1}
\]
We will also use
\[
 D_3=a_1a_2a_3-a_1^2a_4-a_0a_3^2+a_0a_1a_5\ge0.               \tag{2}
\]
These are the second and third leading Hurwitz determinants. The
strict Routh--Hurwitz criterion applies to $P(z+\varepsilon)$ when
all zeros of $P$ lie in the closed left half-plane and $\varepsilon>0$.
Taking $\varepsilon\downarrow0$ gives the necessary weak inequalities.
There is also an elementary proof of (1), the only condition needed
for the main theorem. Normalize to a monic polynomial and put
$D(P)=a_1a_2-a_3$, padding missing coefficients by zeros.
For monic $P,Q$ with next coefficients $u,v$, direct multiplication
gives $D(PQ)=D(P)+D(Q)+uv(u+v)$. A real polynomial with all zeros
in the closed left half-plane factors over the reals into linear and
quadratic monic factors with nonnegative coefficients. Each factor has
$D\ge0$, and the product identity gives (1).

\section{The independent sets near the maximum}
Put $\mu_s=\E X_s=\binom ns q^{\binom{s}{2}}$. For each $k$ we choose
a graph order
\[
 n=n_k=\left\lfloor (k!\lambda_k)^{1/k}q^{-(k-1)/2}\right\rfloor,
 \qquad \lambda_k=\lambda>0 \ \text{or}\ \lambda_k=k .          \tag{3}
\]
In both cases $n\sim(k/e)q^{-(k-1)/2}$. Rounding incurs a relative
error $O(k/n)$ in the $k$th power, and
$(n)_k/n^k=1+O(k^2/n)$ for the falling factorial $(n)_k$. Hence
$\mu_k\sim\lambda_k$. Also
\[
 \frac{\mu_{s+1}}{\mu_s}=\frac{n-s}{s+1}q^s .                  \tag{4}
\]
It follows that $\mu_{k+1}\to0$, even when $\lambda_k=k$, and
$\mu_{k-j}\to\infty$ for every fixed $j\ge1$.

\begin{lemma}\label{lem:counts}
Along (3), with probability tending to one there is no independent
$(k+1)$-set. For every fixed $j\ge 1$,
$X_{k-j}/\mu_{k-j}\to1$ in probability. When $\lambda_k=k$ the same
holds for $j=0$. When $\lambda_k=\lambda$ is constant,
$X_k$ converges in distribution to $\pois(\lambda)$.
\end{lemma}
\begin{proof}
Markov's inequality proves the first assertion. Put $s=k-j$ with
$j\ge0$ fixed. For a fixed $s$-set $S$ and a uniformly random $s$-set
$T$, write $L=|S\cap T|$. The joint probability of independence
of both sets acquires a factor $q^{-\binom L2}$ over the product of the
two probabilities. Consequently
\[
 \frac{\E X_s^2}{\mu_s^2}=\sum_{\ell=0}^s H_\ell,\qquad
 H_\ell=\Prob(L=\ell)q^{-\binom{\ell}{2}}.
\]
The terms for $\ell=0,1$ sum to at most one, and $H_s=1/\mu_s$.
We show that the other terms contribute $o(1)$.

For $2\le\ell\le s/2$, use a union bound over the $\ell$ chosen elements
of $S$. The chance that $T$ contains them is at most $(s/n)^\ell$, so
\[
 H_\ell\le\left(\frac{s^2}{n}q^{-s/4}\right)^\ell .
\]
For $s/2<\ell<s$, write $h=s-\ell$. The hypergeometric count gives
\[
 H_{s-h}
 =\mu_s^{-1}\binom{s}{h}\binom{n-s}{h}
      q^{h(s-(h+1)/2)}
 \le\mu_s^{-1}\left(sn q^{3s/4-1/2}\right)^h .
\]
By (3) the bases in both bounds tend to zero exponentially in $k$
up to polynomial factors, and $\mu_s$ is bounded away from zero.
Summing the two geometric bounds proves the assertion. Thus
$\operatorname{Var}(X_s)/\mu_s^2\le 1/\mu_s+o(1)$, which proves
the concentration claims.

For the Poisson assertion take $s=k$ and $\lambda_k=\lambda$.
For each independent-set indicator use the dependency neighborhood
of $k$-sets meeting its index set in at least two vertices (including
itself). Outside that neighborhood, the entire collection of indicators
depends on edge variables disjoint from those defining this indicator.
The dependency-neighborhood Poisson approximation bounds the total
variation error by a bounded factor times $b_1+b_2$. Here the sum of
products of marginal probabilities is
$b_1=\mu_k^2\Prob(L\ge2)=O(\mu_k^2 k^4/n^2)=o(1)$.
The sum of joint probabilities for distinct neighbors is
$b_2=\mu_k^2\sum_{\ell=2}^{k-1} H_\ell=o(1)$ by the two bounds above.
The approximating mean $\mu_k$ tends to $\lambda$.
[\href{https://doi.org/10.1214/aop/1176991491}{Arratia, Goldstein and Gordon, Theorem 1}]
\end{proof}

\begin{proof}[Proof of Theorem~\ref{thm:main}]
Choose any fixed $0<\lambda<q^2$. Equation (4) gives
\[
 \frac{\mu_{k-1}\mu_{k-2}}{\mu_{k-3}}
 =\mu_k q^{-2}\frac{k(n-k+3)}{(k-2)(n-k+1)}
 \longrightarrow \lambda q^{-2}<1 .
\]
By Lemma~\ref{lem:counts}, with probability tending to one $X_{k-3}>0$
and $X_{k-1}X_{k-2}<X_{k-3}$. On that event, if $X_k\ge1$ and
$X_{k+1}=0$, the polynomial has degree exactly $k$ and
\[
 D_2=X_{k-1}X_{k-2}-X_kX_{k-3}<0 .
\]
Thus $\liminf_{k\to\infty}\Prob(\Unst_{n_k})\ge 1-e^{-\lambda}$.
Taking the supremum over $\lambda<q^2$ proves the stated limit superior.
\end{proof}

Even the qualitative disproof can dispense with Poisson approximation.
Cauchy--Schwarz gives $\Prob(X_k>0)\ge \mu_k^2/\E X_k^2$, whose limit
inferior is at least $\lambda/(1+\lambda)$ by the overlap estimate.
The same argument then gives a lower bound $q^2/(1+q^2)$, exactly
$1/5$ at $p=1/2$.

\section{Two strengthenings}
\begin{corollary}\label{cor:escape}
Fix $0<p<1$. There exist $c_p,\eta_p>0$ and an infinite sequence of
orders such that the probability of a root with
$\Rea z>c_p n/(\log n)^2$ has limit inferior at least $\eta_p$.
\end{corollary}
\begin{proof}
Use (3) with $\lambda=q^2/2$ and put $R=\mu_{k-1}/\mu_k$.
For every fixed $j$, repeated application of (4) gives
\[
 \frac{\mu_{k-j}}{\mu_k R^j}\ \longrightarrow\ q^{\binom j2},
                                                               \tag{5}
\]
since successive backward ratios, divided by $R$, tend in order to
$1,q,q^2,\ldots$. On the event $X_k=1$, $X_{k+1}=0$, set
$Y_j=X_{k-j}/R^j$. Then $Y_0=1$ and $(Y_1,Y_2,Y_3)$ concentrates at
$(\lambda,\lambda q,\lambda q^3)$ even conditionally, since the event
has limiting probability $\lambda e^{-\lambda}>0$.

Translate $I_G$ by $z\mapsto z+\delta R/k$, for fixed $\delta>0$.
Its coefficient of $z^{k-j}$, divided by $R^j$, is
\[
 B_j=\sum_{i=0}^j \binom{k-i}{j-i}(\delta/k)^{j-i}Y_i .
\]
For $j=0,1,2,3$ the limits are polynomials in $\delta$. At $\delta=0$
the limit of $B_1B_2-B_0B_3$ equals $\lambda q(\lambda-q^2)<0$.
It remains strictly negative for some fixed $\delta>0$. This is the
translated second Hurwitz determinant divided by $R^3$, so with
conditional probability tending to one the original polynomial has
a root of real part greater than $\delta R/k$. From (3)--(4),
$R/k\sim e^2 n/k^2$ and $k\sim 2\log n/\log(1/q)$, proving the
claim with $\eta_p=\lambda e^{-\lambda}$.
\end{proof}

\begin{corollary}\label{cor:whp}
If $1-p>\sqrt{(\sqrt5-1)/2}$, then $\limsup_n\Prob(\Unst_n)=1$.
\end{corollary}
\begin{proof}
Take $\lambda_k=k$ in (3). With probability tending to one $\alpha=k$,
and by (5) and Lemma~\ref{lem:counts} the normalized coefficients
$a_j/(\mu_k R^j)$ converge in probability to $q^{\binom j2}$ for
$0\le j\le5$. In (2) all terms scale by the common factor
$\mu_k^3 R^6$. The normalized determinant therefore tends to
\[
 q^4-2q^6+q^{10}=q^4(q^2-1)(q^4+q^2-1),
\]
which is strictly negative under the stated hypothesis.
\end{proof}

\section*{Scope, prior work and reproduction}
The subsequences in (3) suffice to disprove convergence of the stability
probability to one. The full sequence of probabilities at $p=1/2$ is
not determined here. Brown and Nowakowski had proved that almost every
uniform random graph has a \emph{nonreal} independence root and that the
average polynomial has simple real zeros. A nonreal root may still have
negative real part. Chen, Ning and Tu answered the separate complete
bipartite graph question in the same 2018 paper. The threshold and
Poisson arguments here use standard random-clique methods on the
complement graph. The step addressing stability is the Hurwitz test on
the fluctuating \emph{top} coefficient.
[\href{https://doi.org/10.1016/j.jctb.2004.10.001}{Brown and Nowakowski};
\href{https://arxiv.org/abs/2505.24381}{the complete bipartite graph result}]

As an exact finite check, take $p=1/2$, $k=32$, and $\lambda=1/8$
in (3), giving $n=555422$. Put $V_s=\E X_s^2/\mu_s^2$ and
$\varepsilon=1/10$. Direct rational arithmetic in the overlap formula gives
\[
 \frac{(1+\varepsilon)^2\mu_{k-1}\mu_{k-2}}
      {(1-\varepsilon)\mu_{k-3}}<1,\qquad
 \frac1{V_k}-\mu_{k+1}
   -\varepsilon^{-2}\sum_{j=1}^3(V_{k-j}-1)>0.1032 .
\]
The latter expression is a lower bound for $\Prob(\Unst_n)$ by
Cauchy--Schwarz, Markov's inequality, Chebyshev's inequality and a
union bound, by substituting the first displayed bound into (1). The accompanying
\texttt{verify.py} checks both comparisons, both Hurwitz minors, the
coefficient ratios and the overlap moments for every five-vertex graph
at two edge densities. Searches of the source paper,
Cameron's thesis, indexed citing works and publicly searchable research
notes through 18 November 2026 found no earlier resolution of the
random-graph question or sector conjecture. This research note awaits
independent review.
\end{document}
